\documentclass[10pt,a4paper]{amsart}
\usepackage{amsmath,amssymb,amsthm,graphicx}
\usepackage{verbatim}
\newenvironment{reference}{\comment}{\endcomment}
%\newenvironment{reference}{}{}
\newenvironment{slogan}{\comment}{\endcomment}
\newenvironment{history}{\comment}{\endcomment}
% For commutative diagrams we use Xy-pic
\usepackage[all]{xy}

% We use 2cell for 2-commutative diagrams.
\xyoption{2cell}
% Theorem environments.
%
\theoremstyle{plain}
\newtheorem{theorem}[subsection]{Theorem}
\newtheorem{proposition}[subsection]{Proposition}
\newtheorem{lemma}[subsection]{Lemma}

\theoremstyle{definition}
\newtheorem{definition}[subsection]{Definition}
\newtheorem{example}[subsection]{Example}
\newtheorem{exercise}[subsection]{Exercise}
\newtheorem{situation}[subsection]{Situation}

\theoremstyle{remark}
\newtheorem{remark}[subsection]{Remark}
\newtheorem{remarks}[subsection]{Remarks}

\numberwithin{equation}{subsection}

% Macros
%
\def\lim{\mathop{\mathrm{lim}}\nolimits}
\def\colim{\mathop{\mathrm{colim}}\nolimits}
\def\Spec{\mathop{\mathrm{Spec}}}
\def\Hom{\mathop{\mathrm{Hom}}\nolimits}
\def\Ext{\mathop{\mathrm{Ext}}\nolimits}
\def\SheafHom{\mathop{\mathcal{H}\!\mathit{om}}\nolimits}
\def\SheafExt{\mathop{\mathcal{E}\!\mathit{xt}}\nolimits}
\def\Sch{\mathit{Sch}}
\def\Mor{\mathop{\mathrm{Mor}}\nolimits}
\def\Ob{\mathop{\mathrm{Ob}}\nolimits}
\def\Sh{\mathop{\mathit{Sh}}\nolimits}
\def\NL{\mathop{N\!L}\nolimits}
\def\CH{\mathop{\mathrm{CH}}\nolimits}
\def\proetale{{pro\text{-}\acute{e}tale}}
\def\etale{{\acute{e}tale}}
\def\QCoh{\mathit{QCoh}}
\def\Ker{\mathop{\mathrm{Ker}}}
\def\Im{\mathop{\mathrm{Im}}}
\def\Coker{\mathop{\mathrm{Coker}}}
\def\Coim{\mathop{\mathrm{Coim}}}

% Boxtimes
%
\DeclareMathSymbol{\boxtimes}{\mathbin}{AMSa}{"02}

%
% Macros for moduli stacks/spaces
%
\def\QCohstack{\mathcal{QC}\!\mathit{oh}}
\def\Cohstack{\mathcal{C}\!\mathit{oh}}
\def\Spacesstack{\mathcal{S}\!\mathit{paces}}
\def\Quotfunctor{\mathrm{Quot}}
\def\Hilbfunctor{\mathrm{Hilb}}
\def\Curvesstack{\mathcal{C}\!\mathit{urves}}
\def\Polarizedstack{\mathcal{P}\!\mathit{olarized}}
\def\Complexesstack{\mathcal{C}\!\mathit{omplexes}}
% \Pic is the operator that assigns to X its picard group, usage \Pic(X)
% \Picardstack_{X/B} denotes the Picard stack of X over B
% \Picardfunctor_{X/B} denotes the Picard functor of X over B
\def\Pic{\mathop{\mathrm{Pic}}\nolimits}
\def\Picardstack{\mathcal{P}\!\mathit{ic}}
\def\Picardfunctor{\mathrm{Pic}}
\def\Deformationcategory{\mathcal{D}\!\mathit{ef}}
\usepackage{cleveref}
\crefname{originalsection}{original section}{original sections}
\Crefname{originalsection}{Original section}{Original sections}
\makeatletter
\newlabel{flat-section-eta}{{0F8L}{source}{}{sourcepointer-1}{}}
\newlabel{flat-section-eta@cref}{{[originalsection][0][]0F8L}{[1][0][]source}{}{}{}}
\makeatother
\crefname{originaldefinition}{original definition}{original definitions}
\Crefname{originaldefinition}{Original definition}{Original definitions}
\makeatletter
\newlabel{homology-definition-shift-cochain}{{011G}{source}{}{sourcepointer-2}{}}
\newlabel{homology-definition-shift-cochain@cref}{{[originaldefinition][0][]011G}{[1][0][]source}{}{}{}}
\makeatother
\title{The de Rham complex of log poles along an effective Cartier divisor sits in a canonical short exact sequence}
\author{The Stacks Project Authors}
\date{stacks-project:derham.tex:local-snapshot-20261012; GNU FDL 1.2 or later}
\begin{document}
\maketitle
\noindent\textit{Editorial note.} All delivered mathematical text is the author's own, copied from the pinned Stacks Project chapter file derham.tex; no statement or proof step was rewritten, reordered or removed.\par\medskip
\noindent\textit{Editorial note.} Preamble provenance: the pinned chapter file opens with \textbackslash{}input\{preamble\} and carries no preamble text of its own, so the project's single Stacks-wide preamble staged at sources/stacks-derham/preamble.tex is excerpted verbatim as three spans, lines 8-12 (verbatim package and the project's own reference/slogan/history display environments), lines 14-18 (Xy-pic setup) and lines 154-217 (theorem declarations and mathematical macros). The wrapper class and the amsmath/amssymb/amsthm/graphicx base packages are supplied by the assembler. No macro was invented or redefined.\par\medskip
\noindent\textit{Editorial note.} Same-chapter pointers are located at their \textbackslash{}label lines inside the pinned chapter file; cross-chapter pointers are located in the project's own tag list staged in the snapshot (sources/stacks-chow/tags.txt), at the verbatim lines carrying their canonical tags.\par\medskip
\noindent\textit{Allowed standard prerequisites (editorial).}\begin{itemize}
\item The comparison of a module with its twist by a Cartier divisor, giving the inclusion of the de Rham complex. More on Flatness, Section \textbackslash{}ref\{flat-section-eta\} (invoked at proof line 2423). It is the reference for the canonical inclusion \$\textbackslash{}Omega\textasciicircum{}p\_\{X/S\} \textbackslash{}subset \textbackslash{}Omega\textasciicircum{}p\_\{X/S\}(Y)\$ and the statement that the differentials have no nonzero local sections supported on \$Y\$.
\item The shift convention for cochain complexes. Homology, definition homology-definition-shift-cochain (invoked at proof line 2466). It fixes the sign of the differential on \$\textbackslash{}Omega\textasciicircum{}\textbackslash{}bullet\_\{Y/S\}[-1]\$, which the author matches with the forced sign in the Leibniz rule for the log complex.
\item The local splitting hypothesis defining log poles. Same chapter derham.tex, the assumption stated in the statement (used at proof lines 2405-2417). For a local equation \$f\$ it gives \$\textbackslash{}Omega\_\{X/S,y\} = \textbackslash{}mathcal\{O\}\_\{X,y\}\textbackslash{}mathrm\{d\}f \textbackslash{}oplus M\$ with torsion free exterior powers, from which \$\textbackslash{}Omega\textasciicircum{}p\_\{Y/S,y\} = \textbackslash{}wedge\textasciicircum{}p(M/fM)\$ and the absence of sections supported on \$Y\$ follow.
\end{itemize}\medskip
\ifdefined\hypertarget\hypertarget{sourcepointer-1}{}\fi
\noindent\textit{Source reference.} Original section 0F8L: sources/stacks-chow/tags.txt, source lines 19046--19046. Navigation pointer for \textbackslash{}ref\{flat-section-eta\} used inside the counted span (line 2423); the argument that consumes it is retained in the span.\par\medskip
\ifdefined\hypertarget\hypertarget{sourcepointer-2}{}\fi
\noindent\textit{Source reference.} Original definition 011G: sources/stacks-chow/tags.txt, source lines 1362--1362. Navigation pointer for \textbackslash{}ref\{homology-definition-shift-cochain\} used inside the counted span (line 2466); the argument that consumes it is retained in the span.\par\medskip
\begin{lemma}
\label{lemma-log-complex}
Let $X \to S$ be a morphism of schemes. Let $Y \subset X$ be an
effective Cartier divisor.
Assume the de Rham complex of log poles is defined for $Y \subset X$ over $S$.
There is a canonical short exact sequence
of complexes
$$
0 \to \Omega^\bullet_{X/S} \to
\Omega^\bullet_{X/S}(\log Y) \to
\Omega^\bullet_{Y/S}[-1] \to 0
$$
\end{lemma}

\begin{proof}
Our assumption is that for every $y \in Y$ and local equation
$f \in \mathcal{O}_{X, y}$ of $Y$ we have
$$
\Omega_{X/S, y} = \mathcal{O}_{X, y}\text{d}f \oplus M
\quad\text{and}\quad
\Omega^p_{X/S, y} = \wedge^{p - 1}(M)\text{d}f \oplus \wedge^p(M)
$$
for some module $M$ with $f$-torsion free exterior powers $\wedge^p(M)$.
It follows that
$$
\Omega^p_{Y/S, y} = \wedge^p(M/fM) = \wedge^p(M)/f\wedge^p(M)
$$
Below we will tacitly use these facts.
In particular the sheaves $\Omega^p_{X/S}$ have no nonzero local
sections supported on $Y$ and we have a canonical inclusion
$$
\Omega^p_{X/S} \subset \Omega^p_{X/S}(Y)
$$
see More on Flatness, Section \ref{flat-section-eta}. Let $U = \Spec(A)$
be an affine open subscheme such that $Y \cap U = V(f)$ for some
nonzerodivisor $f \in A$. Let us consider the $\mathcal{O}_U$-submodule
of $\Omega^p_{X/S}(Y)|_U$ generated by
$\Omega^p_{X/S}|_U$ and $\text{d}\log(f) \wedge \Omega^{p - 1}_{X/S}$
where $\text{d}\log(f) = f^{-1}\text{d}(f)$.
This is independent of the choice of $f$ as another generator of the
ideal of $Y$ on $U$ is equal to $uf$ for a unit $u \in A$ and we get
$$
\text{d}\log(uf) - \text{d}\log(f) = \text{d}\log(u) = u^{-1}\text{d}u
$$
which is a section of $\Omega_{X/S}$ over $U$. These local
sheaves glue to give a quasi-coherent submodule
$$
\Omega^p_{X/S} \subset \Omega^p_{X/S}(\log Y) \subset \Omega^p_{X/S}(Y)
$$
Let us agree to think of $\Omega^p_{Y/S}$ as a quasi-coherent
$\mathcal{O}_X$-module. There is a unique surjective
$\mathcal{O}_X$-linear map
$$
\text{Res} : \Omega^p_{X/S}(\log Y) \to \Omega^{p - 1}_{Y/S}
$$
defined by the rule
$$
\text{Res}(\eta' + \text{d}\log(f) \wedge \eta) = \eta|_{Y \cap U}
$$
for all opens $U$ as above and all
$\eta' \in \Omega^p_{X/S}(U)$ and $\eta \in \Omega^{p - 1}_{X/S}(U)$.
If a form $\eta$ over $U$ restricts to zero on $Y \cap U$, then
$\eta = \text{d}f \wedge \eta' + f\eta''$ for some forms $\eta'$ and $\eta''$
over $U$. We conclude that
we have a short exact sequence
$$
0 \to \Omega^p_{X/S} \to \Omega^p_{X/S}(\log Y) \to \Omega^{p - 1}_{Y/S} \to 0
$$
for all $p$. We still have to define the differentials
$\Omega^p_{X/S}(\log Y) \to \Omega^{p + 1}_{X/S}(\log Y)$.
On the subsheaf $\Omega^p_{X/S}$ we use the differential of
the de Rham complex of $X$ over $S$. Finally, we define
$\text{d}(\text{d}\log(f) \wedge \eta) = -\text{d}\log(f) \wedge \text{d}\eta$.
The sign is forced on us by the Leibniz rule (on $\Omega^\bullet_{X/S}$)
and it is compatible with the differential on $\Omega^\bullet_{Y/S}[-1]$
which is after all $-\text{d}_{Y/S}$ by our sign convention in
Homology, Definition \ref{homology-definition-shift-cochain}.
In this way we obtain a short exact
sequence of complexes as stated in the lemma.
\end{proof}

\end{document}
